1. Core Definitions
A Riesz space is a real vector space \( E \) equipped with a partial order \( \leq \) that is compatible with the vector structure and forms a lattice.

Compatibility Axioms
For all \( x, y, z \in E \) and \( \lambda \in \mathbb{R}_{\geq 0} \):
- Translation Invariance: If \( x \leq y \), then \( x + z \leq y + z \).
- Positive Homogeneity: If \( x \leq y \), then \( \lambda x \leq \lambda y \).
- Lattice Property: For every pair \( x, y \), the supremum \( x \vee y \) and infimum \( x \wedge y \) exist in \( E \).

The "Big Three" OperationsEvery element \( x \) can be decomposed into its functional components:
- Positive part: \( x^+ = x \vee 0 \)
- Negative part: \( x^- = (-x) \vee 0 \)
- Absolute value: \( |x| = x \vee (-x) \)
- The Jordan Decomposition: \( x = x^+ - x^- \) and \( |x| = x^+ + x^- \).2. 

Orthogonality
Two elements \( x, y \) are disjoint (denoted \( x \perp y \)) if \( |x| \wedge |y| = 0 \).
Disjoint Complement: \( A^d = \{x \in E : |x| \wedge |a| = 0 \text{ for all } a \in A\} \). \( A^d \) is always a band.

Key Substructures:
Understanding how to "slice" a Riesz space is vital for functional analysis.

Riesz Subspace: 
A vector subspace \( L \subseteq E \) such that \( x, y \in L \implies x \vee y \in L \).

Solid Set/Ideal: 
A subset \( S \subseteq E \) s.t. for all \( s \in S \) and  \( x \in E \), \( |x| \leq |s| \implies x \in S \) 
is solid. A solid set which is also a subspace is an ideal.
One can construct the usual vector lattice quotient.

Band: An ideal \( B \) that is "dedekind closed": 
If \( D \subseteq B \) and \( \sup D = x \), then \( x \in B \).
It suffices also if the same holds for \( |x| \).

Notion of "completeness" in Riesz spaces are:
Archimedean: If \( nx \leq y \) for all \( n \in \mathbb{N} \) and \( x \geq 0 \), then \( x = 0 \). (No "infinitely small" elements).
Dedekind Complete: Every non-empty subset bounded above has a supremum.
Dedekind \( \sigma \)-complete: Every countable subset bounded above has a supremum.4. 

Fundamental Theorems

The Riesz Decomposition Property
If \( 0 \leq x \leq y_1 + y_2 \), then there exist \( x_1, x_2 \) such that \( x = x_1 + x_2 \) with \( 0 \leq x_i \leq y_i \). 
Freudenthal’s Spectral Theorem:
In a Dedekind complete Riesz space with a principal unit \( e \), every element \( x \) in the principal ideal generated by \( e \) can be approximated uniformly by "step functions" 
(linear combinations of projections of \( e \) onto bands).

Kakutani’s Representation Theorem
Every AM-space (an \( L^\infty \)-like space) with a unit is lattice-isomorphic to \( C(K) \), the space of continuous functions on a compact Hausdorff space \( K \).

Classification of finite dimensional vector lattices:
Let \( E \) be of dimension \( n \). If \( E \) is archimedian it is isomorphic to \( \mathbb{R}^n \), otherwise
it is isomorphic to some \( \mathbb{R}^k \times \mathbb{R}^{n-k}_L \), where \( \mathbb{R}^{n-k}_L \) has the lexicographical ordering.

Polars and Ideals:
In a locally convex vector lattice \( E \), the polar of every solid neighborhood of the origin is a solid subset of the continuous dual space 
\( E' \). Moreover, the family of all solid equicontinuous subsets of \( X' \) is a fundamental family of equicontinuous sets, 
the polars (in the bidual \( X'' \)) form a neighborhood base of the origin for the natural topology on \( X'' \)
(that is, the topology of uniform convergence on equicontinuous subset of \( X' \).

Common Examples
\( C[0,1] \): Archimedean, but not Dedekind complete (pointwise suprema of continuous functions aren't always continuous).
\( L^p(\mu) \): Dedekind complete for \( 1 \leq p \leq \infty \).
\( \mathbb{R}^n \): With the component-wise ordering.



We now describe some special algebraic properties of subsets of vector spaces. The line segment joining vectors  and  is the set .

Describe some common classification of subsets of vector spaces

A subset \(A\) of a vector space is:
- convex if it includes the line segment joining any pair of its points.
- absorbing (or radial) if for any \(x\) some multiple of \(A\) includes the line segment joining \(x\) and zero.
    That is, if there exists some \(\alpha_0 > 0\) satisfying \(\alpha x \in A\) for every \(0 < \alpha < \alpha_0\).
    Equivalently, \(A\) is absorbing if for each vector \(x\) there exists some \(\alpha_0 > 0\) such that \(\alpha x \in A\) whenever \(-\alpha_0 < \alpha < \alpha_0\).
- circled (or balanced) if for each \(x \in A\) the line segment joining \(x\) and \(-x\) lies in \(A\).
    That is, if for any \(x \in A\) and any \(|\alpha| \leq 1\) we have \(\alpha x \in A\).
- symmetric if \(x \in A\) implies \(-x \in A\).
- star-shaped about zero if it includes the line segment joining each of its points with zero.
    That is, if for any \(x \in A\) and any \(0 \leq \alpha \leq 1\) we have \(\alpha x \in A\).


What is required from a topology on a vector space and what are properties of open/closed subsets in a TVS?

A topology \( \tau \) on a vector space \( X \) is called a linear topology if the operations
addition and scalar multiplication are \( \tau \)-continuous. 
That is, if \( (x, y) \mapsto x + y \), \( (\alpha, x) \mapsto \alpha x \) are continuous.

Most topologies are derived from a norm, that is a functional \( \| \cdot \| : X \to \mathbb{R} \)
satisfying:
1. \( \|x\| \geq 0 \) and equality iff. \( x = 0 \)
2. \( \|\alpha x \| = |\alpha| \|x\| \) for all \( \alpha \in \mathbb{R} \)
2. \( \|x + y \| \leq \|x\| \|y\| \)
The neighborhood basis is then given by the open balls \( B_x(\varepsilon) = \{y \in X : \|x - y\| < \varepsilon \} \).
This always generates a metric with \( d(x, y) = \|x - y\| \) and coincides with it in topology.

In a topological vector space:
1. The algebraic sum of an open set and an arbitrary set is open \(\implies\) translation invariant open sets.
2. Nonzero multiples of open sets are open.
3. If \( B \) is open, then for any set \( A \) we have \( \overline{A} + B = A + B \).
4. The algebraic sum of a compact set and a closed set is closed. (However,
the algebraic sum of two closed sets need not be closed.)
5. The algebraic sum of two compact sets is compact.
6. Scalar multiples of closed sets are closed.
7. Scalar multiples of compact sets are compact.
8. A linear functional is continuous if and only if it is continuous at 0.



State some facts in regards to convex sets in a vector space

A set \( C \subseteq X \) is convex if and only if it contains the line segment for all its points.
Thus for every finite (infinite is not required, except the set is closed for a TVS)
\( \{x_0, \dots, x_n\} \in C \) and \( \{\lambda_0, \dots, \lambda_n\} \subset [0, 1]   \)
with \( \sum \lambda = 1 \), we have \( \sum^{n}_{i=0} \lambda_{i} x_{i} \in C \).

In any vector space:
1. The sum of two convex sets is convex.
2. Scalar multiples of convex sets are convex.
3. A set C is convex if and only if αC + βC = (α + β)C for all nonnegative
scalars α and β.
4. The intersection of an arbitrary family of convex sets is convex.
5. A convex set containing zero is circled if and only if it is symmetric.
6. In a topological vector space, both the interior and the closure of a convex
set are convex.

The smallest convex set containing some set \( S \subseteq X \) is called the convex hull \( \text{ch}(S) \) of \( S \).
(Similar notions for circled and closed).

The convex hull and the convex circled hull of a compact subset
of a finite dimensional vector space are compact sets.
For inifinite dimensional spaces we require more:
In a completely metrizable locally convex space, 
the closed convex hull of a compact set is compact.

If \( S \) is finite the resulting \( \text{ch}(S) \) is called a polytope.
Every polytope in a tvs is compact.

Carathéodory’s Convexity Theorem in an \( n \)-dimensional vector space:
Every vector in the convex hull of a nonempty set can be written as a convex
combination using no more than \( n+1 \) vectors from the set.



What is a bounded set in a topological vector space?

A subset \( A \) of a topological vector space \( (X, \tau) \) is (topologically)
bounded, or more specifically \( \tau \)-bounded, if for each neighborhood \( V \) of zero
there exists some \( \lambda > 0 \) such that \( A \subset \lambda V \).

This can also be characterized with sequences:
A set \( B \) is bounded if and only if for every sequence \(\{x_n\} \subseteq B\) 
and every sequence of scalars \(\{\alpha_n\} \to 0\), the sequence \(\{\alpha_n x_n\} \to 0\) in \( X \).

Every subset of a bounded set is bounded.
The union of a finite number of bounded sets is bounded.
If \( A \) and \( B \) are bounded, then \(A + B = \{a + b : a \in A, b \in B\}\) is bounded.
The closure \(\overline{B}\) of a bounded set \( B \) is also bounded.
Every compact set in a TVS is bounded. 

Heine-Borel
In a TVS \( X \), every closed and bounded set is compact if and only if \( X \) is finite-dimensional.

Kolmogorov’s Normability Criterion
A Hausdorff topological vector space \( X \) is normable 
(its topology can be generated by a norm) if and only if it 
possesses a convex, bounded neighborhood of zero.

Uniform boundedness principle
If \( X \) is a F-space (complete metrizable TVS), \( Y \) a TVS and \( \mathcal{F} \subseteq \mathcal{L}(X, Y) \)
then if \( \{ T x : T \in \mathcal{F}\}  \) is bounded \( Y \) for all \( x \in X \), then \( \mathcal{F} \) is equicontinuous.

If the topology of \( X \) is induced by a metric \( d \), every "TVS-bounded" set is "metric-bounded" \(\forall x, y \in B, d(x, y) < M \). 
However, the converse is not always true unless the metric satisfies a set of additional conditions.

The generalization of boundedness to non-normable TVS is referred to as Bornology.


